\documentclass[UTF8,12pt,a4paper,fontset=fandol]{ctexart}

% 支持从项目根目录、深度学习目录或当前目录编译。
\makeatletter
\def\input@path{{../}{../../}}
\makeatother
\usepackage{researchnotes}

\title{模型剪枝的原理与前置知识}
\author{}
\date{}

\begin{document}
\maketitle

\begin{abstract}
模型剪枝通过删除参数、连接或计算结构，在预测质量与资源成本之间寻求可接受的折中。理解剪枝不能只记忆“小权重可以删除”：参数大小依赖网络的表示方式，连接的重要性依赖数据分布，局部扰动不等于恢复训练后的损失，而数学上的稀疏也不自动带来硬件加速。本文从具体线性模型出发，介绍线性代数、概率统计、微分与优化、神经网络计算图及计算机系统等前置知识，推导幅值剪枝、一阶敏感度、二阶损失近似、误差补偿、稀疏正则化和层级预算分配，继而讨论结构化剪枝、初始化剪枝、动态稀疏训练以及大语言模型的训练后剪枝。文中给出可核查的算例、完整算法流程、部署与实验原则、常见误区和练习解答，区分精确结论、局部近似、算法设计与经验性假设。
\end{abstract}
\tableofcontents

\section{研究对象、基本问题与学习路线}
\label{sec:prune-introduction}

\subsection{从删除一个系数开始}
考虑线性预测器 $f(\mathbf x)=2x_1+0.01x_2$。如果两个输入分量的典型绝对值均约为一，把第二个系数改为零，预测通常只发生很小的变化；但如果 $x_2$ 的典型绝对值约为一千，同样的操作就可能删掉数量级为十的预测贡献。因此，参数小只是一个容易计算的信号，并不是重要性小的定义。更进一步，若 $x_2=x_1$ 恒成立，原函数也可写成 $2.01x_1$，删除第二个连接后适当修改第一个系数，就能完全恢复原有预测。这里同时出现了剪枝的三个核心因素：权重大小、输入统计以及剩余参数的补偿能力。

\keyterm{模型剪枝（model pruning）}是按照某种规则移除模型中的连接、参数组或计算模块，并使所得模型满足目标质量与资源约束的方法集合。连接被移除时，可以先用二值掩码把相应权重固定为零；通道或模块被移除时，还需要重建相邻张量和计算图。前者改变有效参数，后者可以进一步改变算子的实际形状。是否需要重新训练、能否得到更小的文件、是否降低真实延迟，取决于剪枝粒度、保存格式和执行环境，不能由“删除了一半参数”单独决定。

剪枝通常不是寻找唯一正确的子网络，而是在给定数据、训练预算和部署条件下选择一个合适的子网络。即使两个子网络具有相同参数量，也可能具有不同的表达能力、优化难度和运行时间。一个面向手机图像分类的方案与一个面向服务器文本生成的方案，资源瓶颈和质量指标都可能不同。因而研究问题应先写成“在什么任务和硬件上，保留什么质量，减少哪一种成本”，再讨论具体重要性分数。

\subsection{与其他压缩方法的关系}
\keyterm{量化（quantization）}减少权重或激活所需的数值精度，剪枝减少被保留的元素或结构；\keyterm{低秩分解（low-rank factorization）}用较低秩矩阵的乘积近似原矩阵，剪枝则约束某些坐标或参数组为零。一个稠密的秩一矩阵可以没有任何零元素，一个对角矩阵可以很稀疏却具有满秩，因此低秩与稀疏是不同结构。\keyterm{知识蒸馏（knowledge distillation）}利用教师输出或表示指导学生训练，主要改变学习信号，可以在剪枝后的恢复训练中使用。

\keyterm{随机失活（dropout）}在训练时随机屏蔽部分激活以形成正则化，常规推理通常仍使用完整网络。固定剪枝掩码则永久限制某些连接参与计算。\keyterm{条件计算（conditional computation）}根据输入选择部分模块，例如混合专家模型的路由；它的活跃结构可能随样本变化。静态剪枝固定部署结构，动态路由选择当前路径，两者可以组合，但参数总量、每次活跃参数量与计算量应分别统计。

\subsection{全篇的知识依赖}
本文先把神经网络写成带参数的函数，明确损失、梯度、曲率、掩码和预算，然后从“删除后发生多大变化”推导重要性判据，再从“其余参数怎样补偿”引入二阶方法和局部重建。随后讨论如何删除完整结构、怎样在训练过程中改变稀疏模式，以及如何将这些思想应用于语言模型。最后把数学压缩与存储、访存、稀疏算子和实验评估联系起来。读者可以先掌握幅值、一阶与局部重建三条主线，再阅读二阶推导和稀疏训练的进阶部分。

\section{线性代数与网络表示的前置知识}
\label{sec:prune-linear}

\subsection{向量、矩阵、张量和掩码}
本文将单个输入写为列向量 $\mathbf x\in\mathbb R^d$。线性层为 $\mathbf y=\mathbf W\mathbf x+\mathbf b$，其中 $\mathbf W\in\mathbb R^{m\times d}$、$\mathbf b\in\mathbb R^m$，输出属于 $\mathbb R^m$。矩阵的第 $i$ 行决定第 $i$ 个输出，第 $j$ 列决定第 $j$ 个输入对所有输出的影响。删除一个权重只删去一个标量连接；删除一行改变输出维度；删除一列改变输入依赖。这三种操作对后续层的影响完全不同。

\keyterm{张量（tensor）}在本文中指带有多个轴的数值数组。例如卷积权重可写为 $\mathbf W\in\mathbb R^{C_{\rm out}\times C_{\rm in}\times k_h\times k_w}$，依次对应输出通道、输入通道与卷积核的两个空间维度。图像批量激活常写成 $B\times C\times H\times W$，其中 $B$ 是批量大小。实际框架可能采用其他轴顺序，所以任何“按第几个维度剪枝”的实现，都应先核对该维度的语义，而不是凭数组位置推断。

设模型共有 $p$ 个待考虑的标量参数，把它们展平为 $\theta\in\mathbb R^p$。\keyterm{掩码（mask）}为 $\mathbf m\in\{0,1\}^p$，有效参数是 $\widetilde\theta=\mathbf m\odot\theta$，其中 $\odot$ 表示逐元素乘法。$m_j=0$ 表示第 $j$ 个连接不参与模型计算。用 $\operatorname{supp}(\mathbf m)=\{j:m_j=1\}$ 表示保留的支持集。掩码保留某个位置并不要求该位置当前数值一定非零，因而掩码密度与数值非零比例在严格意义上可能不同。

\subsection{范数、稀疏度和误差的度量}
向量的 $\ell_1$ 范数为 $\|\theta\|_1=\sum_j|\theta_j|$，$\ell_2$ 范数为 $\|\theta\|_2=(\sum_j\theta_j^2)^{1/2}$。常用记号 $\|\theta\|_0=\#\{j:\theta_j\ne0\}$ 表示非零元素个数，但它不满足范数的齐次性，严格说不是范数。矩阵的 Frobenius 范数为 $\|\mathbf A\|_F=(\sum_{i,j}A_{ij}^2)^{1/2}$；谱范数为 $\|\mathbf A\|_2=\sup_{\|\mathbf x\|_2=1}\|\mathbf A\mathbf x\|_2$，用于控制矩阵对输入误差的最大放大倍数。

相对于选定参数集合的\keyterm{稀疏率（sparsity）}定义为 $s=1-\|\mathbf m\|_0/p$，保留比例或密度为 $\rho=1-s$。若嵌入层、偏置和归一化参数不参与剪枝，分母是否包含它们会改变报告结果。局部层有百分之九十的稀疏率，不代表全模型也有百分之九十的稀疏率；应同时说明可剪参数的数量、实际保留数量以及完整模型的参数总量。

参数误差、层输出误差和任务误差属于不同空间。$\|\widetilde\theta-\theta\|_2$ 小，不必然使所有输入上的预测误差都小；$\|\widetilde{\mathbf W}\mathbf X-\mathbf W\mathbf X\|_F$ 小，只能直接说明给定输入矩阵上的线性输出接近；分类准确率还取决于输出之间的相对大小。选择范数时必须同时说明衡量的对象，不能笼统称为“模型距离”。

\subsection{相关性、冗余与可替代性}
令随机输入为 $\mathbf X$，假设 $\mathbb E\|\mathbf X\|_2^2<\infty$，定义二阶矩矩阵 $\mathbf C=\mathbb E[\mathbf X\mathbf X^\top]$。它是半正定矩阵，因为任意 $\mathbf v$ 满足 $\mathbf v^\top\mathbf C\mathbf v=\mathbb E[(\mathbf v^\top\mathbf X)^2]\geq0$。若输入均值为零，$\mathbf C$ 就是协方差矩阵；否则二阶矩还包含均值外积。线性输出重建通常直接涉及二阶矩，不能未经说明就把它替换成中心化协方差。

对线性预测器 $f(\mathbf X)=\mathbf w^\top\mathbf X$，参数变化 $\delta$ 引起的平均平方输出误差恰好为
\begin{equation}
 \mathbb E[(\delta^\top\mathbf X)^2]=\delta^\top\mathbf C\delta.
 \label{eq:prune-linear-error}
\end{equation}
若只删除第 $j$ 个参数，$\delta=-w_j\mathbf e_j$，误差为 $w_j^2C_{jj}$；若同时删除两个参数，还会出现 $2w_iw_jC_{ij}$。其中 $\mathbf e_j$ 是第 $j$ 个标准基向量。交叉项既可能增大误差，也可能减小误差，说明“分别不重要”不能简单推出“同时删除仍不重要”。

\begin{example}[相关特征允许补偿]
若 $X_2=X_1$ 且 $\mathbb E[X_1^2]=1$，模型为 $3X_1+2X_2$。仅把第二个权重置零，平均平方输出误差为 $4$；若同时把第一个权重改成 $5$，输出误差变成零。若改为 $X_1,X_2$ 不相关、均值为零且方差均为一，无论如何修改第一个权重，遗漏 $2X_2$ 所产生的最小均方误差仍为 $4$。可补偿程度由数据中的线性依赖决定，不能仅由参数数量判断。
\end{example}

\section{概率、学习目标与数据划分}
\label{sec:prune-learning}

\subsection{总体风险与经验风险}
设输入与标签的随机对为 $(\mathbf X,Y)\sim\mathcal D$，模型是 $f_\theta$，损失函数为 $\ell(f_\theta(\mathbf x),y)$。\keyterm{总体风险（population risk）}为 $\mathcal R(\theta)=\mathbb E_{(\mathbf X,Y)\sim\mathcal D}[\ell(f_\theta(\mathbf X),Y)]$，前提是期望存在。给定 $n$ 个样本组成的训练集 $\{(\mathbf x_a,y_a)\}_{a=1}^n$，\keyterm{经验风险（empirical risk）}为
\begin{equation}
 L(\theta)=\frac1n\sum_{a=1}^n\ell(f_\theta(\mathbf x_a),y_a).
 \label{eq:prune-empirical}
\end{equation}
剪枝算法通常只能访问有限样本，所以它估计的是与当前样本相关的重要性。训练样本上损失不变，不能推出部署分布上的损失也不变。

分类模型常输出实数向量 $\mathbf z\in\mathbb R^K$，称为 logits，经 softmax 得到 $p_\theta(k\mid\mathbf x)=\exp(z_k)/\sum_{r=1}^K\exp(z_r)$。真实标签为 $y$ 时，交叉熵损失是 $-\log p_\theta(y\mid\mathbf x)$。准确率仅检查最大概率对应的类别是否正确，而交叉熵还关注概率大小。剪枝可能不改变预测类别，却降低正确类别的置信程度，因此只看准确率会忽略部分变化；反过来，少数边界样本上的小幅输出变化也可能造成准确率下降。

\subsection{训练、校准、验证和测试承担不同职责}
\keyterm{校准数据（calibration data）}在剪枝语境中通常用于收集激活、梯度、曲率或局部重建目标，不一定包含人工标签。幅值方法可以不使用校准数据，激活加权方法需要输入，监督损失敏感度通常还需要标签或教师信号。校准数据不等同于测试数据；使用测试集来选择剪枝率、阻尼系数或恢复轮数，会让最终测试结果偏向已经被选择过程反复利用的样本。

一个清晰的流程是：训练集用于参数学习，校准子集用于构造剪枝统计量，验证集用于选择预算和超参数，测试集在方案固定后评估。校准子集可以从训练集抽取，但应记录抽样方式与样本量；若它与恢复训练数据重叠，也应明确说明。对于语言模型，还应记录词元数、长度分布、语言和任务来源。使用相同数量的短句与长文档，所得到的上下文覆盖和有效词元数量并不相同。

\subsection{估计误差与分布偏移}
设连接 $j$ 的单样本重要性是 $A_j(\mathbf X,Y)$，用 $n_c$ 个独立同分布校准样本的均值 $\widehat a_j$ 估计其期望。如果方差有限，则 $\operatorname{Var}(\widehat a_j)=\operatorname{Var}(A_j)/n_c$。这个等式只说明固定连接估计量的方差如何随样本量变化，并不保证对成千上万个连接排序时所有名次都稳定。分数接近的连接尤其容易因采样波动而互换位置，必要时应比较多个校准子集所得的掩码和质量。

语言序列中的词元并非独立样本，图像增强产生的多个视图也存在相关性，所以不能直接把词元数或视图数当作独立观测数代入上述方差关系。部署分布若与校准分布不同，增加同一狭窄来源的样本也未必解决问题。一个仅在常见类别上活跃的参数，与一个只对罕见但重要样本起作用的参数，可能具有完全不同的平均分数与最坏情况作用；剪枝目标应与实际质量约束相匹配。

\section{微分、曲率与优化的前置知识}
\label{sec:prune-calculus}

\subsection{梯度与反向传播}
对可微标量函数 $L:\mathbb R^p\to\mathbb R$，梯度 $\mathbf g=\nabla L(\theta)$ 的第 $j$ 个分量是 $\partial L/\partial\theta_j$。小扰动 $\delta$ 产生的一阶变化约为 $\mathbf g^\top\delta$。梯度下降以 $\theta_{t+1}=\theta_t-\eta_t\mathbf g_t$ 更新参数，$\eta_t>0$ 为学习率。\keyterm{反向传播（backpropagation）}是沿计算图应用链式法则计算梯度的算法；它提供局部导数，不会直接回答删除一大组参数并重新训练后的最优损失。

若输出为 $\mathbf z(\theta)\in\mathbb R^K$，Jacobian 矩阵 $\mathbf J=\partial\mathbf z/\partial\theta\in\mathbb R^{K\times p}$，则 $\nabla_\theta L=\mathbf J^\top\nabla_{\mathbf z}L$。当剪枝对象是某个中间通道时，可以为通道引入一个标量门，再对这个门求导；这样就能把结构变化转化为普通标量敏感度。门的位置必须固定，因为在线性变换之前、归一化之后或残差相加之前插入门，所表示的删除操作不同。

\subsection{Taylor 展开与 Hessian 矩阵}
若 $L$ 在 $\theta$ 的邻域内二阶连续可微，Hessian 矩阵 $\mathbf H=\nabla^2L(\theta)$ 的元素为二阶偏导数，并有
\begin{equation}
 L(\theta+\delta)-L(\theta)
 =\mathbf g^\top\delta+\frac12\delta^\top\mathbf H\delta
 +o(\|\delta\|_2^2).
 \label{eq:prune-taylor}
\end{equation}
若 Hessian 在连接 $\theta$ 与 $\theta+\delta$ 的区域内以常数 $M$ 满足 Lipschitz 条件，余项的绝对值还可界为 $M\|\delta\|_2^3/6$。删除权重对应 $\delta_j=-\theta_j$，它未必是小扰动，因此使用二阶评分时仍需用实际损失检验近似是否可靠。

在局部极小点，若可微则梯度为零，若二阶可微则 Hessian 半正定；但训练终点可能仅是随机梯度较小的点，也可能受到正则项、早停和数值误差影响。ReLU 网络的损失通常只有分片光滑性，扰动跨越激活状态边界时，固定区域内的 Taylor 模型不再直接适用。文中的二阶推导在指定光滑性与正定性条件下成立，将它应用于实际网络是一种近似建模。

\subsection{Hessian、Gauss--Newton 与 Fisher 的区别}
对单样本复合损失 $\ell(\mathbf z(\theta),y)$，二阶链式法则给出
\begin{equation}
 \nabla_\theta^2\ell
 =\mathbf J^\top\nabla_{\mathbf z}^2\ell\,\mathbf J
 +\sum_{k=1}^K\frac{\partial\ell}{\partial z_k}\nabla_\theta^2z_k.
 \label{eq:prune-hessian-decomposition}
\end{equation}
对样本平均第一项得到\keyterm{广义 Gauss--Newton 矩阵（generalized Gauss--Newton matrix）}。如果损失关于输出 $\mathbf z$ 是凸的，该矩阵半正定，但它忽略了输出关于参数的二阶导数项，通常不等于完整 Hessian。局部线性模型或某些特殊情形下两者可以重合，不能把这一性质无条件推广到整个深度网络。

对条件概率模型，\keyterm{Fisher 信息矩阵（Fisher information matrix）}可写为 $\mathbf F=\mathbb E_{\mathbf X}\mathbb E_{Y\sim p_\theta(\cdot\mid\mathbf X)}[\mathbf u\mathbf u^\top]$，其中 $\mathbf u=\nabla_\theta\log p_\theta(Y\mid\mathbf X)$。在支持集不随参数变化、可交换积分与微分等正则条件下，负对数似然的 Hessian 在模型分布下取期望，结果等于 Fisher。用实际标签梯度外积构造的 $\widehat{\mathbf F}=n^{-1}\sum_a\mathbf g_a\mathbf g_a^\top$ 常称经验 Fisher；它半正定，却一般既不等于完整 Hessian，也不等于模型 Fisher。

这些替代矩阵的价值在于较稳定、较便宜或可用矩阵向量积计算，而不是把不同数学对象变成同一个对象。特别是 $\bigl(n^{-1}\sum_a g_{aj}\bigr)^2$ 与 $n^{-1}\sum_a g_{aj}^2$ 不相等：前者允许样本梯度相互抵消，后者衡量单样本梯度的平方大小。二者会产生不同的剪枝排序，论文和代码必须明确所用的是哪一种统计量。

\section{把剪枝写成一个优化问题}
\label{sec:prune-optimization}

\subsection{固定预算、惩罚目标与恢复训练}
在最简单的非结构化情形，给定保留数量 $k\leq p$，剪枝问题可以写为
\begin{equation}
 \min_{\theta,\,\mathbf m\in\{0,1\}^p}L(\mathbf m\odot\theta)
 \quad\text{满足}\quad\|\mathbf m\|_0\leq k.
 \label{eq:prune-budget}
\end{equation}
这里同时选择结构和参数。固定已经训练好的 $\theta$ 只优化 $\mathbf m$，得到的是不恢复参数的剪枝问题；先选择掩码再在该掩码下更新权重，则近似求解内层连续优化。理论上的 $\min_\theta$ 与有限步恢复训练也不同，后者依赖初始化、优化器和训练预算。

更一般地，令 $\mathcal M$ 为满足计算图与硬件限制的掩码集合，$C(\mathbf m)$ 为成本，可用 $\min L(\mathbf m\odot\theta)$ 且 $C(\mathbf m)\leq B_0$ 表达资源约束。$B_0$ 是资源预算。也可以优化 $L(\mathbf m\odot\theta)+\lambda C(\mathbf m)$，但在离散、非凸问题中，改变 $\lambda$ 未必能获得每一个指定预算下的最优解，惩罚形式与约束形式不能无条件互换。实测延迟还依赖张量形状与运行条件，通常不是每个连接成本的简单相加。

\subsection{三种不同的“重要性”}
固定其他参数时，删除集合 $P$ 的真实即时损失变化是 $L(\theta-\theta^{[P]})-L(\theta)$，这里 $\theta^{[P]}\in\mathbb R^p$ 表示只在 $P$ 上保留原参数值、其余位置为零的向量。允许充分优化其余参数时，关心的则是约束最优值与原模型损失之差。实际算法为了节省计算，往往使用幅值、梯度、重建误差等代理分数。这三者分别对应即时影响、可恢复影响和可计算估计，不应混称为同一个精确量。

重要性还依赖删去的其他集合。某个连接独立删除几乎无影响，可能是因为另一个连接能替代它；两个都删去时影响却很大。因此一般不能给每个参数赋予一个永久不变、与上下文无关的重要性。逐步剪枝、每轮重新评分和参数补偿，都是在有限成本下适应这种相互依赖；它们并不保证找到全局最佳支持集。

\subsection{全局、逐层与组内选择}
\keyterm{全局剪枝（global pruning）}把多个层的候选参数放在一起排序，让不同层获得不同稀疏率；\keyterm{逐层剪枝（layer-wise pruning）}先给每层指定预算，再在层内排序。全局方法较灵活，但分数尺度不同会偏向某些层，也可能把关键瓶颈层几乎清空。逐层方法更容易维持最低容量，却可能把不敏感层的富余预算留住，同时过度压缩敏感层。

\keyterm{组内剪枝（group-wise pruning）}还可以要求每个输出神经元保留固定数量的输入，或每个连续小块保留固定数量的元素。这类约束缩小了搜索空间，可能更利于计算，也限制了算法将预算集中到少数位置的能力。比较方法时应匹配可行集合：同样百分之五十稀疏率，任意位置删除与每四个位置固定保留两个，不是同一道优化问题。

\section{幅值剪枝与最简单的可靠基线}
\label{sec:prune-magnitude}

\subsection{幅值排序究竟优化了什么}
\keyterm{幅值剪枝（magnitude pruning）}按照 $a_j=|\theta_j|$ 排序，保留分数最大的若干权重。它不需要数据或反向传播，评分成本与参数量成正比；全排序通常需要 $O(p\log p)$ 次比较，而选择阈值可以使用更高效的选择算法。幅值法的数学依据并非“训练过的小权重必然没用”，而是它精确解决了一个参数空间中的最近稀疏近似问题。

\begin{proposition}[欧氏距离下的最佳稀疏近似]
\label{prop:prune-hard-threshold}
给定 $\theta\in\mathbb R^p$ 与整数 $0\leq k\leq p$，问题 $\min_{\mathbf v:\|\mathbf v\|_0\leq k}\|\theta-\mathbf v\|_2^2$ 的一个解是保留 $|\theta_j|$ 最大的 $k$ 个坐标，并令其余坐标为零。若阈值处有相同幅值，最优支持集可以不唯一。
\end{proposition}
\begin{proof}
固定支持集 $S$ 时，各坐标可独立优化。$j\in S$ 上令 $v_j=\theta_j$，误差最小；$j\notin S$ 上只能取零，剩余误差为 $\sum_{j\notin S}\theta_j^2$。在至多保留 $k$ 项的条件下，应使被保留的平方和最大，即选择幅值最大的坐标。该结论没有使用任务损失，因此也没有保证任务损失最小。
\end{proof}

如果进一步假设局部损失为 $L(\theta+\delta)-L(\theta)\approx c\|\delta\|_2^2/2$，其中所有方向的曲率都相同且 $c>0$，则幅值法也近似最小化局部损失。这是强假设。不同层的输入尺度、不同方向的曲率和参数间的相关性，都会破坏这种各向同性。实践中幅值法仍值得作为基线，因为它简单、容易审计，并能检验复杂方法的额外统计成本是否确实带来了收益。

\subsection{重参数化反例}
设一个两层、无偏置的网络为 $f(\mathbf x)=\mathbf W_2\sigma(\mathbf W_1\mathbf x)$，激活函数为 ReLU，满足 $\sigma(a\mathbf u)=a\sigma(\mathbf u)$，其中 $a>0$。把某个隐藏神经元对应的第一层行乘以 $a$，第二层列除以 $a$，网络函数保持不变。但当 $a$ 很小时，第一层这一行的权重会变得很小，跨层幅值排序就更可能删除它。函数完全相同的两个参数表示，因尺度不同而得到不同剪枝结果，说明原始幅值不是函数层面的不变量。

可以采用逐层预算、组归一化或结合输入激活的评分缓解尺度问题，但每一种处理都引入新的选择。除以层内权重范数会改变层间可比性，却不一定对应某个真实损失。对包含归一化的网络，缩放对输出的影响还可能被归一化部分抵消。解释归一化评分时应说明它解决的是哪种尺度不匹配，而不要把处理后的数值称为绝对真实的重要性。

\subsection{迭代幅值剪枝}
一个实用基线是先训练稠密模型，再删除当前较小的一部分有效权重，在剩余支持集上恢复训练，重复直至达到目标稀疏率。这类训练、剪枝、再训练流程的代表工作包括 Han 等人的连接剪枝研究\cite{han2015}。逐步过程让剩余参数有机会吸收误差，也让后续排名基于更新后的参数。它的成本包括初始训练与每轮恢复，不能只计算最后稀疏模型的一次训练成本。

若每轮删除当前剩余参数的比例 $r$，执行 $T$ 轮后密度为 $(1-r)^T$，稀疏率为 $1-(1-r)^T$。例如每轮删除百分之二十，五轮后保留比例为 $0.8^5=0.32768$，稀疏率为 $67.232\%$，并不是百分之一百。实现中应按目标累计保留数选择支持集，明确取整规则与并列分数处理，否则多个轮次之后实际稀疏率可能与计划不一致。

\section{一阶敏感度、激活统计与门控评分}
\label{sec:prune-first-order}

\subsection{从删除扰动到权重乘梯度}
删除权重 $\theta_j$ 对应 $\delta=-\theta_j\mathbf e_j$。由式\eqref{eq:prune-taylor}的一阶项得
\begin{equation}
 \Delta L_j\approx-\theta_jg_j.
 \label{eq:prune-first-score}
\end{equation}
若希望识别“删除可能引起较大改变”的位置，可以使用 $|\theta_jg_j|$；若目标是直接估计带符号的损失增加，则应保留 $-\theta_jg_j$ 的符号。这两种目标不同：绝对值会把可能降低损失的删除也视为重要，而带符号估计在局部近似不稳定时可能过于乐观。

给第 $j$ 个参数引入连续门 $c_j$，有效参数为 $c_j\theta_j$，在 $\mathbf c=\mathbf 1$ 处有 $\partial L(\mathbf c\odot\theta)/\partial c_j=\theta_jg_j$。因此权重乘梯度也是损失对连接门的敏感度。SNIP 在初始化附近使用连接敏感度选择稀疏结构\cite{lee2019}；它要选择的是后续可以训练的支持集，而不是保护一个已训练模型的全部当前功能。相同公式在不同时间点使用，算法目标与解释都不相同。

\subsection{为什么训练终点的一阶评分可能退化}
在精确驻点，完整训练目标的梯度为零，一阶删除预测全部为零，但删除参数通常仍会增加损失。这不是反向传播失效，而是一阶展开本来就无法描述驻点附近的主要二阶变化。实际算法可能采用单样本梯度绝对值的平均、平方梯度或小批量统计，以避免跨样本抵消。它们引入了不同代理目标，不能继续声称等于完整经验损失的一阶变化。

例如两个样本对同一参数的梯度分别为 $2$ 与 $-2$，平均梯度为零，平均绝对梯度为 $2$，平均平方梯度为 $4$。如果先对批量平均再取绝对值，结果为零；若先逐样本取绝对值再平均，结果为 $2$。批量大小、损失的求和或平均方式以及词元掩码，都会影响分数尺度。复现实验时应记录这些细节，尤其不能把不同批量归约得到的分数直接混合比较。

\subsection{对通道和特征使用 Taylor 评分}
设一个通道的激活向量为 $\mathbf h_c$，给整个通道乘标量门 $z_c$，其余计算保持不变。将 $z_c$ 从一改为零时，一阶变化为 $-\sum_u h_{c,u}\,\partial L/\partial h_{c,u}$，其中 $u$ 遍历该通道的空间位置或词元位置。可以先在单样本内求和，再跨样本取绝对值或平方平均；也可以先取逐位置绝对值，但这又是不同分数。基于 Taylor 展开的通道重要性是卷积网络剪枝的一条代表路线\cite{molchanov2017}。

仅以平均激活大小作为分数，可以发现长时间不活跃的通道，却看不到下游对该通道的敏感度。一个激活很大但被下一层乘以极小系数的通道，影响可能有限；一个激活较小但下游放大很多倍的通道，影响可能很大。激活与梯度的乘积同时包含前向信号与局部反向敏感度，但仍依赖当前数据和当前网络状态。

\subsection{训练轨迹与移动剪枝}
预训练模型在新任务上微调时，参数当前幅值反映了历史训练的积累，未必反映其对新任务的作用。\keyterm{移动剪枝（movement pruning）}通过微调过程中的一阶信息学习掩码分数\cite{sanh2020}。为理解其动机，考察梯度下降小步更新：$\Delta(\theta_j^2/2)\approx-\eta\theta_jg_j$。若 $-\theta_jg_j$ 为正，当前更新倾向于增大权重的绝对值，说明任务梯度正在推动该连接远离零。

这一解释只是一阶、小步长下的局部分析。真实算法学习连续分数并用离散选择形成掩码，常借助直通估计器传播近似梯度；它不等于简单比较微调前后的绝对值差，也不保证所有远离零的参数都应保留。适用背景、分数更新和稀疏调度共同构成方法，不能仅凭一个“变化量”就认定实现了原算法。

\section{二阶剪枝：曲率、补偿与最优性边界}
\label{sec:prune-second-order}

\subsection{对角曲率与 Optimal Brain Damage}
假设当前点近似驻点，忽略梯度项和高阶余项，再用 Hessian 的对角部分近似完整曲率，删除第 $j$ 个参数的预测代价为
\begin{equation}
 a_j^{\rm diag}=\frac12H_{jj}\theta_j^2.
 \label{eq:prune-obd}
\end{equation}
这一思想对应 Optimal Brain Damage，简称 OBD\cite{lecun1989}。当所有 $H_{jj}$ 相同且为正时，它退化为按幅值排序；当曲率不同时，小权重也可能位于敏感方向。对角近似忽略参数相互作用，尤其不能表达另一个权重变化后对当前删除的补偿。

作为自设算例，设两个参数为 $\theta_1=0.1$、$\theta_2=1$，局部 Hessian 为 $\operatorname{diag}(1000,1)$。幅值法先删除第一个参数，而二阶估计对应的损失增加分别为 $5$ 和 $0.5$，于是曲率法先删除第二个参数。这里的结论对给定二次模型是精确的；对一般网络，它仅说明幅值和曲率可能给出不同排序，不能据此断言曲率法在所有任务中更优。

\subsection{允许其他参数调整的 Optimal Brain Surgeon}
现在保留完整的二次模型，并允许所有未被约束的参数同步调整。设 $\mathbf H$ 对称正定、梯度为零，删除第 $j$ 个参数的局部问题是
\begin{equation}
 \min_\delta\frac12\delta^\top\mathbf H\delta
 \quad\text{满足}\quad\mathbf e_j^\top\delta=-\theta_j.
 \label{eq:prune-obs-problem}
\end{equation}
\begin{proposition}[单坐标删除的最优二次补偿]
\label{prop:prune-obs}
在上述假设下，令 $q_j=(\mathbf H^{-1})_{jj}>0$，则唯一最优解及其代价为
\begin{equation}
 \delta^*=-\frac{\theta_j}{q_j}\mathbf H^{-1}\mathbf e_j,
 \qquad a_j^{\rm OBS}=\frac{\theta_j^2}{2q_j}.
 \label{eq:prune-obs}
\end{equation}
\end{proposition}
\begin{proof}
引入乘子 $\nu$，Lagrange 函数为 $\tfrac12\delta^\top\mathbf H\delta+\nu(\mathbf e_j^\top\delta+\theta_j)$。驻点条件给出 $\delta=-\nu\mathbf H^{-1}\mathbf e_j$，代入约束得 $\nu=\theta_j/q_j$。再代回二次目标得式\eqref{eq:prune-obs}。由于 $\mathbf H$ 正定，目标严格凸且约束为非空仿射集合，所以该驻点是唯一全局最优解。
\end{proof}

这一局部补偿思想对应 Optimal Brain Surgeon，简称 OBS\cite{hassibi1992}。名称中的最优是针对所建立的二次约束问题，不是整个深度网络剪枝问题的全局最优。真实网络可能有奇异或不定 Hessian，且同时删除很多参数会远离展开点；实现通常需要替代曲率、阻尼和近似更新。

\subsection{一个可以逐项核对的补偿例子}
设 $\theta=(1,1)^\top$，$\mathbf H=\left(\begin{smallmatrix}2&1\\1&2\end{smallmatrix}\right)$，则 $\mathbf H^{-1}=\tfrac13\left(\begin{smallmatrix}2&-1\\-1&2\end{smallmatrix}\right)$。删除第一个参数但不补偿时 $\delta=(-1,0)^\top$，二次损失增加为 $1$。允许第二个参数补偿时，式\eqref{eq:prune-obs}给出 $\delta^*=(-1,1/2)^\top$，代价为 $3/4$。第二个权重从 $1$ 变成 $3/2$，通过曲率交叉项减小了损失增加。

若同时删除两个参数，则只能取 $\delta=(-1,-1)^\top$，总代价为 $3$；而两个单独删除的 OBS 代价之和为 $3/2$。差异来自两次单独删除都假设另一个参数仍能补偿，同时删除时这一条件不再成立。因而把所有单坐标代价一次求出后直接相加，只是一个排序近似，不能当成联合删除的精确损失。

\subsection{多个约束与组剪枝的推广}
设要删除的坐标集合为 $P$，$\mathbf E_P$ 由这些标准基向量作为列组成，$\theta_P=\mathbf E_P^\top\theta$ 为对应参数子向量。相同的乘子推导得到
\begin{align}
 \delta^*&=-\mathbf H^{-1}\mathbf E_P
 (\mathbf E_P^\top\mathbf H^{-1}\mathbf E_P)^{-1}\theta_P,\label{eq:prune-group-update}\\
 a_P&=\frac12\theta_P^\top
 (\mathbf E_P^\top\mathbf H^{-1}\mathbf E_P)^{-1}\theta_P.
 \label{eq:prune-group-cost}
\end{align}
正定性保证中间矩阵可逆。这个公式揭示了组内相关性的作用，但枚举所有候选组仍可能非常昂贵，尤其完整 Hessian 的存储需要 $O(p^2)$ 空间。块对角近似、低秩近似和矩阵向量积方法可以降低成本，同时也改变了所使用的曲率模型。

\subsection{非驻点、阻尼与数值稳定性}
如果梯度不为零，应保留 $\mathbf g^\top\delta$。对正定 $\mathbf H$，令 $\mathbf d=-\mathbf H^{-1}\mathbf g$，单坐标约束问题的解为 $\delta^*=\mathbf d-\frac{\theta_j+d_j}{q_j}\mathbf H^{-1}\mathbf e_j$。此时存在一个不带删除约束的下降步骤，不能再把全部目标变化都解释为剪枝的代价。使用零梯度 OBS 公式之前，应检查研究对象是已经拟合的局部重建问题，还是尚在变化的全局训练损失。

常见稳定化方式是以 $\mathbf H_\lambda=\mathbf H+\lambda\mathbf I$ 替代曲率，$\lambda>0$ 为阻尼系数。若原矩阵半正定，这保证正定；若原矩阵不定，$\lambda$ 必须大于最负特征值绝对值才能保证正定。阻尼同时惩罚过大的补偿步，因此改变了优化目标。数值实现宜通过 Cholesky 分解或线性方程求解获得所需结果，而不是机械地显式求逆；遇到病态矩阵时，还应检查输入统计、尺度和精度。

\section{局部重建与输入激活加权}
\label{sec:prune-reconstruction}

\subsection{把一层剪枝转化为稀疏回归}
对线性层 $\mathbf W\in\mathbb R^{m\times d}$，将 $n_c$ 个校准输入作为列，构成 $\mathbf X_c\in\mathbb R^{d\times n_c}$。原层输出为 $\mathbf Y_c=\mathbf W\mathbf X_c$，目标是找稀疏矩阵 $\widehat{\mathbf W}$ 使
\begin{equation}
 E(\widehat{\mathbf W})=\frac{1}{2n_c}
 \|\widehat{\mathbf W}\mathbf X_c-\mathbf W\mathbf X_c\|_F^2
 \label{eq:prune-reconstruct}
\end{equation}
较小。该目标不需要真实标签，且关于每一行权重都是精确二次函数。令 $\mathbf C_c=\mathbf X_c\mathbf X_c^\top/n_c$，某一行变化为列向量 $\delta$ 时误差是 $\delta^\top\mathbf C_c\delta/2$，因此局部 Hessian 正是输入二阶矩，不能把它称作整个语言模型任务损失的 Hessian。

这种重建将复杂网络中的问题简化为局部输出匹配，代价是忽略下游对不同输出方向的敏感程度。若后续映射对某个方向特别敏感，该方向很小的误差也可能有较大的任务影响；另一些方向的误差可能被后续归一化或投影削弱。因此局部重建通常是可计算的代理目标，而非端到端任务损失的精确等价形式。

\subsection{固定支持集后的精确重建}
考虑一行原始输出 $\mathbf y=\mathbf X_c^\top\mathbf w$，保留输入坐标集合为 $S$。令 $\mathbf X_{c,S}$ 表示只保留 $S$ 中行的输入矩阵，待求权重是 $\mathbf v\in\mathbb R^{|S|}$。若 $\mathbf X_{c,S}\mathbf X_{c,S}^\top$ 可逆，最小二乘解满足
\begin{equation}
 \mathbf v^*=(\mathbf X_{c,S}\mathbf X_{c,S}^\top)^{-1}
 \mathbf X_{c,S}\mathbf y.
 \label{eq:prune-least-squares}
\end{equation}
这是对 $\|\mathbf X_{c,S}^\top\mathbf v-\mathbf y\|_2^2$ 求导并令梯度为零得到的正规方程解。不可逆时仍存在最小二乘解，但可能不唯一；可用伪逆选取最小范数解，或加入正则化获得唯一解。

注意，给定支持集后求权重是连续最小二乘问题，寻找最佳支持集则是组合选择问题。前者有闭式表达式，并不意味着后者也容易求解。直接选择一个掩码后做回归，与在每次选择后重新估计误差并补偿，是不同算法；当每一行选择的支持集不同，所需解的子矩阵也不同，计算成本随之上升。

\subsection{Wanda 分数的基本机制}
只删除单个权重 $W_{ij}$、固定其他权重时，式\eqref{eq:prune-reconstruct}增加量精确等于 $W_{ij}^2\|\mathbf X_{c,j,:}\|_2^2/(2n_c)$。开平方并去掉统一常数得到分数
\begin{equation}
 a_{ij}=|W_{ij}|\,\|\mathbf X_{c,j,:}\|_2.
 \label{eq:prune-wanda}
\end{equation}
Wanda 将权重幅值与相应输入激活范数相乘，并在每个输出维度内选择保留权重\cite{sun2024}。因此一个小权重如果连接到幅度很大的输入，也可能获得较高分数。该思路无需对保留权重进行恢复更新，但同时删除多个连接时，相关输入之间的交叉项仍未被独立排序准确描述。

作为独立算例，设两个权重幅值分别为 $0.02$ 与 $0.5$，相应校准输入范数为 $100$ 与 $1$。幅值法优先删除第一个，激活加权分数分别为 $2$ 和 $0.5$，则优先删除第二个。这个例子只说明输入尺度会改变排序。若校准数据没有包含某个重要输入模式，再精确的范数计算也无法补充没有被观测到的行为。

\subsection{SparseGPT 与逐步误差补偿}
SparseGPT 从层输出重建出发，利用输入曲率信息和近似二阶补偿，在处理矩阵列的过程中选择稀疏位置并修正尚未固定的权重\cite{frantar2023}。其关键是共享和复用曲率计算，配合列顺序处理、块更新与稳定的矩阵分解，使较大矩阵的重建可计算。它不是对全模型完整 Hessian 执行精确 OBS，也不是对所有可能掩码求全局最优。

理解其实现时，应区分逆矩阵对角元素与逆矩阵的 Cholesky 因子对角元素：对当前剩余变量的精确单坐标二次问题，代价含 $1/(\mathbf H^{-1})_{jj}$；将逆曲率序列信息编码进三角因子后，代码里可能出现相应因子对角元素的平方。两种写法依赖不同对象，不能看到平方就直接判断公式矛盾。无梯度再训练也不等于不修改权重：解析或近似误差补偿本身就会改变保留权重。

\subsection{逐层校准与误差传播}
剪完前一层后，后一层接收的输入分布已经改变。可以始终使用原模型激活，强调对原始表示的匹配；也可以顺序运行剪枝后的前缀，用更新后的输入统计拟合当前层。若希望当前层逼近原网络的原始输出，输入和目标可能来自不同前缀，这又构成不同的回归问题。实际流程必须明确输入由哪个模型生成、目标由哪个模型生成，并保持校准与部署时的归一化及随机层模式一致。

假设后续子网络 $g$ 在所考虑的表示区域内是 $K_g$-Lipschitz，即 $\|g(\mathbf u)-g(\mathbf v)\|_2\leq K_g\|\mathbf u-\mathbf v\|_2$，则当前层输出扰动至多被放大 $K_g$ 倍。多层连续剪枝还会累积多个传播项。深层网络的最坏情况常数可能非常松，不能据此精确预测准确率，但它说明了为什么只检查每层误差平均值仍不足以替代端到端验证。

\section{通过正则化学习稀疏结构}
\label{sec:prune-regularization}

\subsection{从约束到稀疏惩罚}
与训练后才挑选删除位置不同，稀疏学习在训练目标中加入偏好，例如 $L(\theta)+\lambda\|\theta\|_1$。$\lambda\geq0$ 控制任务拟合与参数压缩的权衡。$\ell_1$ 惩罚在零点具有尖角，可以通过次梯度或近端方法使坐标精确为零；普通 $\ell_2$ 惩罚则通常使参数连续缩小而不精确归零。二者都可能改变参数分布，但“权重衰减”和“得到一个固定数量的零”不是同一件事。

为了看到差别，考虑一维子问题 $\min_v\tfrac12(v-u)^2+\lambda|v|$。$v>0$ 时驻点为 $u-\lambda$，要求 $u>\lambda$；$v<0$ 时驻点为 $u+\lambda$，要求 $u<-\lambda$；$v=0$ 时次梯度条件为 $|u|\leq\lambda$。于是得到\keyterm{软阈值（soft thresholding）}
\begin{equation}
 \operatorname{soft}_{\lambda}(u)
 =\operatorname{sign}(u)\max(|u|-\lambda,0).
 \label{eq:prune-soft}
\end{equation}
其中 $\operatorname{sign}(0)=0$。对 $\ell_2$ 惩罚 $\lambda v^2/2$，相应解为 $u/(1+\lambda)$，除非 $u=0$，否则通常不会得到零。

\subsection{近端梯度与硬阈值的区别}
\keyterm{近端梯度（proximal gradient）}先对光滑部分走一步 $\mathbf u_t=\theta_t-\eta_t\nabla L(\theta_t)$，再逐坐标应用 $\operatorname{soft}_{\eta_t\lambda}$。它与“普通梯度训练很多步后按幅值剪一次”不同，稀疏惩罚持续参与优化。在光滑凸目标且步长满足相应 Lipschitz 限制等标准条件下，可以建立全局收敛结论；深度网络一般非凸，不能直接继承凸问题的全局最优保证。

若约束是 $\|\theta\|_0\leq k$，一种迭代方案是梯度更新后保留幅值最大的 $k$ 个坐标，即\keyterm{迭代硬阈值（iterative hard thresholding）}。命题\ref{prop:prune-hard-threshold}说明这一步是欧氏意义上的稀疏投影，但投影集合非凸，支持集也可能在迭代间改变。稀疏回归中某些额外条件可以支持理论保证，任意神经网络不自动满足这些条件。

\subsection{组稀疏正则化}
如果部署希望删除完整通道，就应把同一通道的参数视为一个组。对不重叠分组 $\mathcal G$，可使用\keyterm{组稀疏惩罚（group sparsity penalty）} $\lambda\sum_{G\in\mathcal G}\alpha_G\|\theta_G\|_2$，其中 $\alpha_G>0$ 为组权重。组内 $\ell_2$ 范数允许多个参数协同，组间的求和则鼓励整个组归零。结构稀疏学习的一条代表路线是将这种思想用于网络中的通道、滤波器等结构\cite{wen2016}。

对单组近端子问题 $\tfrac12\|\mathbf v-\mathbf u\|_2^2+\lambda\alpha_G\|\mathbf v\|_2$，若 $\mathbf u\ne0$，最优解为 $(1-\lambda\alpha_G/\|\mathbf u\|_2)_+\mathbf u$；若 $\mathbf u=0$，解也为零。这里 $(a)_+=\max(a,0)$。其原因是固定 $\|\mathbf v\|_2$ 时，与 $\mathbf u$ 同方向使平方距离最小，剩下的问题就退化为非负半轴上的一维软阈值。重叠组不能直接逐组套用这个简单公式，因为同一参数会同时受到多个组的约束。

\subsection{归一化尺度与可学习门}
批量归一化（batch normalization, BN）的单通道形式为 $y_c=\gamma_c(x_c-\mu_c)/\sqrt{\sigma_c^2+\varepsilon}+\beta_c$。Network Slimming 对通道尺度 $\gamma_c$ 引入稀疏约束，再据此选择通道\cite{liu2017}。但 $\gamma_c=0$ 时输出是常数 $\beta_c$，不一定为零。真正删除通道需要处理常数项、后续非线性和相邻算子；仅把尺度改成零，不能在所有结构中等价为删除该通道。剪枝后 BN 运行统计也可能需要在合适数据上重新估计。

另一条路线是为参数组加入随机门 $z_G$，优化 $\mathbb E_z[L(\theta,z)]+\lambda\sum_G\Pr(z_G\ne0)$。期望非零组数可以表达结构成本，但离散 Bernoulli 门不方便直接使用路径导数。基于 hard concrete 等连续松弛并截断的门控方法，为零提供非零概率，同时允许学习分布参数\cite{louizos2018}。训练时的随机门、推理时的确定门以及最终结构导出应分别定义；软门很小只表示贡献被缩放，尚不等于算子已经被物理删除。

\section{结构化剪枝与计算图依赖}
\label{sec:prune-structured}

\subsection{非结构化、结构化与半结构化}
\keyterm{非结构化剪枝（unstructured pruning）}独立删除标量权重，通常具有较大的选择自由，但不规则索引增加存储与计算难度。\keyterm{结构化剪枝（structured pruning）}删除整个神经元、通道、注意力头或模块，使输出网络可以由更小的稠密算子构成。\keyterm{半结构化剪枝（semi-structured pruning）}在规则小组内约束零模式，位于二者之间。块稀疏则按矩形块或其他预定块组织非零位置，可减少索引数量并提高计算规律性。

本文约定 $N{:}M$ 表示在指定分组轴上每 $M$ 个元素保留 $N$ 个位置，因而名义稀疏率为 $1-N/M$。实际文献有时用 $N$ 表示被删数量，阅读时必须核对定义；硬件还可能要求特定轴、矩阵形状、对齐和数值类型。即使两个掩码都含一半零，也不意味着都满足相同的半结构化算子要求。

\subsection{线性层和卷积层的联动删除}
对于 $\mathbf h=\sigma(\mathbf W_1\mathbf x+\mathbf b_1)$、$\mathbf y=\mathbf W_2\mathbf h+\mathbf b_2$，其中 $\mathbf W_1\in\mathbb R^{h\times d}$、$\mathbf W_2\in\mathbb R^{m\times h}$，删除隐藏索引集合 $P$ 时，应同时删除 $\mathbf W_1$ 的对应行、$\mathbf b_1$ 的对应元素和 $\mathbf W_2$ 的对应列。得到的模型内部宽度变小，但输入、输出接口不变。这是一个简单的依赖闭包：一个结构变化必须连同所有直接消费该结构的位置一起处理。

卷积中删除输出通道，还会影响下一层的输入通道、当前偏置与相应归一化参数。普通卷积的单样本乘加次数约为 $H_{\rm out}W_{\rm out}C_{\rm out}C_{\rm in}k_hk_w$，忽略偏置与其他算子。减少输出通道不仅节省当前层计算，也可能节省后续层计算；但对分组卷积、深度可分离卷积，通道分组与逐通道对应关系需要额外保持，不能照搬普通卷积的任意切片规则。

\subsection{残差、拼接与归一化}
残差相加 $\mathbf y=\mathbf x+F(\mathbf x)$ 要求两个分支的形状与通道语义对应。只删除 $F$ 输出的某些通道，跳连保持原宽度，通常会产生形状不匹配。可以保持残差主干宽度，仅剪内部扩展层；若要剪主干宽度，则必须沿相连分支传播一致索引。拼接操作虽然不要求分支宽度相同，但会改变后续输入索引的偏移，导出时也要同步更新。

层归一化（layer normalization, LN）沿特征维度计算均值和方差。删除特征后，保留下来的特征也会使用不同的归一化统计，因此不能认为“仅删掉了被选中的坐标”。RMS 归一化同样依赖整个归一化维度的平方均值。结构化剪枝所涉及的不只是参数数组切片，还有算子的定义与统计范围；一个形状合法的网络，也可能与预期的掩码网络具有不同函数。

\subsection{组评分、成本与图约束}
通道评分可以来自权重范数、门梯度、激活重建或组二阶代价；组的大小不相同时，直接比较范数会受元素数量影响。除以组大小或按计算节省归一化，实际上是在改变选择偏好。成本还存在交互：相邻层同时缩窄后，总乘加节省不能总由每组独立节省相加得到。实践中应在候选结构生成后重新计算完整模型成本，避免重复计算共享收益。

完整导出过程需要从保留索引重建权重、偏置、归一化状态以及模型配置，并检查权重共享和所有分支的依赖。若存在一个严格等价的零门参考模型，应在相同输入上比较导出前后的输出；对于改变归一化维度等不保持严格等价的操作，应改为明确评估误差，而不能要求本来就不成立的逐元素相等。

\section{剪枝时机、恢复训练与稀疏调度}
\label{sec:prune-training}

\subsection{训练前、训练中与训练后}
按时间区分，训练前剪枝先确定一个稀疏子网络，再从初始权重开始训练；训练中剪枝使稀疏率或拓扑随训练变化；训练后剪枝压缩已经学好的模型，可能只需校准，也可能还进行恢复训练。\keyterm{一次性剪枝（one-shot pruning）}通常指在一个主要压缩阶段达到目标稀疏率，但不一定意味着算法内部只执行一次排序，更不一定意味着完全不访问数据。不同论文对术语使用有所差异，应以实际流程判断。

\keyterm{恢复训练（recovery training）}或微调使保留参数在新结构下重新适应任务。它需要回答是否继承原权重、是否重置优化器状态、学习率从何处开始以及训练多少步。较低学习率有利于保持原函数附近的行为，但不一定足以适应大幅结构变化；较高学习率可能探索新的解，也可能破坏已保留的能力。这些是需要验证的设计选择，而不是由剪枝定义自动决定的规则。

\subsection{固定掩码为什么需要持续施加}
对有效参数 $\widetilde\theta=\mathbf m\odot\theta$，固定掩码下链式法则给出 $\partial L/\partial\theta_j=m_j\,\partial L/\partial\widetilde\theta_j$。掩码为零的位置接收不到任务损失梯度，但底层存储值仍可能受到原有动量、优化器状态或其他更新项影响。只要每次前向都施加掩码，这些值不会参与输出；如果中途去掉掩码，则可能重新出现有效连接。

另一种实现直接保持参数为零，在更新前屏蔽梯度，并在更新后重新投影到固定支持集。使用带动量或自适应状态的优化器时，还应处理对应位置的状态，尤其在拓扑变化与重新激活连接时明确状态是清零还是继承。PyTorch 的剪枝接口将原参数与掩码组合成重参数化；移除这种重参数化可以把当时的零写回参数，但不会自动缩小稠密张量的形状\cite{pytorchprune}。此后若继续普通训练而不再施加约束，零值通常也不具有永久不变性。

\subsection{渐进稀疏率与预算控制}
设开始剪枝的步数为 $t_0$，结束步数为 $t_1>t_0$，起始与目标稀疏率为 $s_0\leq s_f$。一种可用的平滑日程为
\begin{equation}
 s(t)=s_f+(s_0-s_f)
 \left(1-\frac{t-t_0}{t_1-t_0}\right)^3,
 \qquad t_0\leq t\leq t_1.
 \label{eq:prune-schedule}
\end{equation}
区间之前取 $s_0$，之后取 $s_f$。它是工程调度选择，不是收敛定理。较早阶段给予模型适应时间，结束后留下固定结构的恢复阶段；更新掩码的频率则在统计计算开销与结构适应速度之间折中。

实际执行时按 $\lfloor(1-s(t))p\rfloor$ 或事先规定的其他规则确定累计保留数，之后只在允许的候选集合内选择。若要求掩码单调减少，应保证新支持集包含于旧支持集；若允许再生，则属于不同训练方案。已经为零的参数、受保护层、必须保留的通道数和硬件对齐约束，都应参与预算核算，避免最终统计与计划脱节。

\subsection{蒸馏恢复与附加模块的成本}
设原模型为教师，其类别分布是 $q_T(k\mid\mathbf x)$，剪枝学生分布为 $p_T(k\mid\mathbf x)$，$T>0$ 是 softmax 温度。一个恢复目标可写为
\begin{equation}
 L_{\rm rec}=L_{\rm task}+\alpha T^2
 \mathbb E_{\mathbf X}\!\left[
 \sum_k q_T(k\mid\mathbf X)
 \log\frac{q_T(k\mid\mathbf X)}{p_T(k\mid\mathbf X)}\right],
 \label{eq:prune-distill}
\end{equation}
其中 $\alpha\geq0$ 控制教师约束。教师固定，学生学习保留原函数的重要行为；真实标签和教师信号不一致时，两项目标可能竞争。温度因子是常用的梯度尺度调整方式，并不保证不同温度下完全相同的优化过程。

也可以添加低秩适配模块来补偿剪枝误差。若最终有效权重为 $\mathbf m\odot\mathbf W+\mathbf B\mathbf A$，稠密低秩项可能使原来为零的位置重新变成非零。把它直接合并到基础矩阵后，原稀疏格式可能不再适用；若作为额外分支执行，则有额外计算与内存。评价时应以最终部署结构统计成本，并说明恢复模块是否被保留、合并或再次约束。

\section{彩票假说、初始化剪枝与动态稀疏训练}
\label{sec:prune-sparse-training}

\subsection{彩票假说讨论的是哪一种可训练性}
\keyterm{彩票假说（lottery ticket hypothesis）}提出，在某些随机初始化的稠密网络中，可以找到配合原始初始化使用的稀疏子网络，使其独立训练后达到与原网络可比的性能\cite{frankle2019}。这里的“票”既包括支持集，也包括对应的初始权重。把训练完的大网络剪枝后继续微调，检验的是当前权重的可恢复性；把保留权重重置到初始值再训练，检验的是稀疏子网络的可训练性，两个实验不能相互替代。

一种寻找方式是反复训练、按幅值剪枝，再把保留权重重置到记录的初始化。也可以把权重回退到较早的训练检查点，称为权重回溯；把学习率恢复到早期日程而保留当前权重，则是另一操作。寻找子网络的过程可能耗费多轮训练，所以发现可训练子网络本身并不证明总训练成本已经降低。假说也不是关于所有网络、任务、稀疏率与优化器的普遍定理。

\subsection{初始化剪枝关心学习信号能否传播}
初始化阶段的网络尚未学好任务，此时删除后损失较小不代表最终可训练性一定好。SNIP 使用连接敏感度，是一种任务相关的局部选择；另一条思路是保留训练梯度的传播。令 $G(\theta)=\tfrac12\|\nabla L(\theta)\|_2^2$，在二阶可微时 $\nabla G=\mathbf H\mathbf g$，故删除第 $j$ 个连接的局部变化约为 $-\theta_j(\mathbf H\mathbf g)_j$。这一计算说明了梯度保持评分为何会涉及 Hessian 向量积，而不必显式储存整个 Hessian。

GraSP 从保留梯度流的角度设计初始化剪枝\cite{wang2020}。梯度大有时意味着可学习信号充足，但也可能意味着不稳定；梯度范数只是可训练性的代理指标。若一次全局排序把某个必要层全部删除，输入到输出的路径可能断开，即发生层崩塌。仅优化总体分数而不检查连接结构，在很高稀疏率下尤其容易出现这种问题。

\subsection{不使用任务数据的路径评分}
在简化的无偏置前馈网络中，可以把权重取绝对值、输入设为全一向量，并考察所有路径权重乘积之和。对串联矩阵，这个量写为 $R=\mathbf 1^\top|\mathbf W_L|\cdots|\mathbf W_1|\mathbf 1$，每一项对应一条输入到输出路径的非负贡献。分数 $\theta_j\,\partial R/\partial\theta_j$ 描述某条连接对总路径量的贡献；因为使用绝对值，普通正负权重之间的抵消被移除。

SynFlow 通过这样的路径传播思想和迭代稀疏化避免过早清空网络层\cite{tanaka2020}。上述矩阵表达式用于解释简单网络中的机制，不直接覆盖带归一化、复杂分支或所有非线性的实际架构。无需任务数据意味着评分阶段不读取真实样本，并不意味着稀疏模型此后无需训练，也不意味着评分已经知道哪些路径最符合真实任务。

\subsection{动态稀疏训练与连接再生}
\keyterm{动态稀疏训练（dynamic sparse training）}在保持大致固定活跃参数预算的同时，周期性删除一部分连接并激活新的连接。RigL 是一条代表路线：利用权重幅值删除弱连接，再利用梯度信息选择生长位置\cite{evci2020}。直观上，当前小权重可能贡献有限，而一个未连接位置若具有较大潜在梯度，激活它之后可能快速降低损失。支持集可以变化，使优化不必永久困在一开始选择的稀疏结构中。

这里必须区分底层掩码参数的梯度与有效稠密权重空间中的梯度。固定零掩码会使前者为零，不能直接拿它给所有未连接位置评分。生长步骤需要专门计算或估计这些位置的潜在梯度，其成本与普通稀疏训练步不同。还要确定再生权重初值、优化器状态、拓扑更新间隔和停止再生时刻。若实现仍执行稠密前后向，只是在结果上乘掩码，理论上的活跃乘加减少可能尚未变成实际训练加速。

\section{Transformer 与大语言模型中的剪枝}
\label{sec:prune-transformer}

\subsection{前馈层的宽度剪枝}
\keyterm{变换器（Transformer）}包含注意力模块、前馈模块、归一化和残差连接。设主干隐藏宽度为 $d$，前馈中间宽度为 $h$，普通前馈映射可写为 $\mathbf y=\mathbf W_{\rm down}\sigma(\mathbf W_{\rm up}\mathbf x)$，其中 $\mathbf W_{\rm up}\in\mathbb R^{h\times d}$、$\mathbf W_{\rm down}\in\mathbb R^{d\times h}$。删除中间神经元时，联动删除上投影的行和下投影的列即可保留外部宽度 $d$。偏置存在时也要同步切片。

门控前馈层可以写成 $\mathbf y=\mathbf W_{\rm down}[\sigma(\mathbf W_{\rm gate}\mathbf x)\odot(\mathbf W_{\rm up}\mathbf x)]$，其中两个输入投影均为 $h\times d$。此时中间坐标是两路分量的乘积，删除索引集合 $P$ 必须同时处理门投影与上投影的行，以及下投影的列。只删一路不仅会改变维度对应，也可能不等价于删除预期神经元。评分可以对乘积后的中间激活设置门，从而直接衡量这个结构单元的作用。

\subsection{注意力头、投影矩阵与接口}
为描述序列，令隐藏矩阵 $\mathbf Z\in\mathbb R^{T_s\times d}$ 的每行对应一个词元，$T_s$ 是序列长度。单头注意力使用 $\mathbf Q=\mathbf Z\mathbf W_Q$、$\mathbf K=\mathbf Z\mathbf W_K$、$\mathbf V=\mathbf Z\mathbf W_V$，其中三个投影矩阵属于 $\mathbb R^{d\times d_h}$。忽略掩码书写，其输出为 $\operatorname{softmax}(\mathbf Q\mathbf K^\top/\sqrt{d_h})\mathbf V$；因果语言模型还需在 softmax 之前屏蔽未来位置。多头输出拼接后经输出投影回到主干宽度。

在标准多头结构中删除一个完整头，通常需要同时删除查询、键和值投影对应的头子块，并删除输出投影中读取该头的子块。保持主干宽度不变时，只改变注意力内部宽度。实际实现可能把多个投影融合成一个张量，轴方向也可能与本文列向量约定相反，应按语义定位切片，而不是照抄行列编号。改变头内部维度时，还涉及缩放因子、位置编码的成对维度和后端对齐要求。

\subsection{分组查询、键值共享与缓存}
在\keyterm{分组查询注意力（grouped-query attention, GQA）}中，多个查询头共享一组键和值。查询头与键值头不再一一对应，所以删除某个查询头未必删除相应的缓存；删除一个键值头则可能影响整个查询头组。设批量为 $B$、缓存长度为 $T_s$、层数为 $L_b$、键值头数为 $n_{\rm kv}$、每头维度为 $d_h$、每元素字节数为 $b$，不计元数据时，键和值缓存约需 $2BL_bT_sn_{\rm kv}d_hb$ 字节。

对投影矩阵做非结构化权重剪枝，通常保持输出形状不变，因此不会自动减小上述缓存公式。前馈宽度剪枝也不直接改变注意力缓存。若目标瓶颈是长上下文的缓存内存，应检查剪枝是否确实改变 $n_{\rm kv}$、$d_h$、$L_b$ 或缓存精度，而不能用权重非零数量代替缓存大小。缓存压缩与权重剪枝可以联合研究，但它们约束的是不同对象。

\subsection{层剪枝、残差旁路与深度}
若一个模块采用 $\mathbf h_{l+1}=\mathbf h_l+F_l(\mathbf h_l)$，删除残差分支后可以保留恒等旁路，但模块内部与外部归一化的位置会影响实际变换。前置归一化（pre-norm）与后置归一化（post-norm）结构不能用同一条替换规则概括：这里的前后指归一化相对残差子层的位置，与训练或推理阶段无关。相邻层表示相似可作为候选信号，却不证明该层可无损删除，因为这一层可能专门改变少数重要方向，或只在少见输入上发挥作用。

删除完整层往往能改变串行计算深度，但恢复难度可能较大。可先在验证数据上进行单层删除试验，再评估组合删除和恢复训练；组合结果仍不能由单层影响简单相加。层编号、缓存组织和模型配置必须同步修改，尤其在生成模型中，要确认后续逐词元解码使用的状态与导出结构一致。

\subsection{语言模型质量的多维评估}
给定测试序列中的 $N_t$ 个有效目标词元，平均负对数似然为 $\operatorname{NLL}=-(1/N_t)\sum_t\log p(x_t\mid x_{<t})$，\keyterm{困惑度（perplexity）}为 $\exp(\operatorname{NLL})$。困惑度比较必须匹配分词器、上下文长度、滑动窗口和有效词元计数方式；不同分词器对应的每词元困惑度通常不能直接横向比较。填充位置不应参与平均，文档边界也不能无意引入原任务不存在的上下文。

困惑度衡量给定语料上的预测质量，不完整代表指令遵循、推理、代码、长上下文或罕见知识等能力。某一平均值变化很小，仍可能掩盖特定任务的大幅退化。校准与评估至少应覆盖预期使用场景，并记录生成温度、输出长度等设置。对于混合专家模型，删除专家还需要处理路由概率、负载和专家覆盖；某个专家平均使用频率低，不等于它负责的少数输入可以忽略。

\section{从稀疏模型到真实存储与推理成本}
\label{sec:prune-systems}

\subsection{参数量、文件大小与运行内存}
稠密矩阵有 $mn$ 个元素，若每元素用 $b_v$ 字节，权重存储约为 $mnb_v$。把其中很多值写成零而保持原数组格式，字节数通常不变。若保存为行压缩稀疏格式（compressed sparse row, CSR），需要保存非零值、列索引与行指针。记非零数为 $q$，每索引用 $b_i$ 字节，则简化存储量为
\begin{equation}
 B_{\rm CSR}\approx q(b_v+b_i)+(m+1)b_i.
 \label{eq:prune-csr}
\end{equation}
块稀疏或规则半结构化格式的元数据不同，不能套用同一个公式。

例如一个 $1000\times1000$ 矩阵，用每值两字节的稠密表示需要 $2{,}000{,}000$ 字节。若密度为 $0.2$，每列索引与行指针均为四字节，CSR 约需 $200{,}000\times6+1001\times4=1{,}204{,}004$ 字节，只缩小约 $1.66$ 倍，远小于非零数量减少的五倍。忽略行指针时，这种格式必须满足 $\rho<b_v/(b_v+b_i)=1/3$ 才比稠密更省存储。这里使用十进制字节数，不混用 MB 与 MiB。

模型文件大小也不等于运行峰值内存。推理可能还需要激活、缓存、稀疏格式转换副本和算子工作区；训练还需要梯度、主权重和优化器状态。一个稀疏检查点很小，却在加载时转回稠密数组，就可能没有运行内存收益。应分别测量序列化文件、加载后的常驻内存与任务执行峰值。

\subsection{乘加次数与实际延迟}
\keyterm{乘加次数（multiply--accumulate operations, MACs）}与浮点操作次数的统计口径应声明：有的统计把一次乘加视为两个 FLOPs，有的工具采用不同约定。稀疏矩阵若有 $q$ 个非零元素，与一个向量相乘的理想乘加数约为 $q$；但稠密内核仍可能执行所有零对应的操作。非零数减少只是具备减少理论工作量的可能性，能否实现还取决于选择了怎样的内核。

粗略地，计算耗时受到计算量除以有效吞吐率、数据量除以有效带宽这两个下界的共同限制，还叠加调度与同步开销。小批量解码可能更受权重和缓存访存影响，大批量预填充可能更充分利用矩阵计算；这是一种需通过具体形状测量的分析框架，而非对所有硬件的固定结论。不规则稀疏会引入索引读取、负载不均、非连续访存与较低利用率，过小的矩阵还可能被启动开销主导。

\subsection{Amdahl 定律解释端到端上限}
设原系统耗时中比例 $a$ 的部分可以被剪枝加速，其余比例 $1-a$ 不变，被加速部分获得 $r$ 倍速度，则总加速比为
\begin{equation}
 S=\frac{1}{(1-a)+a/r}.
 \label{eq:prune-amdahl}
\end{equation}
例如 $a=0.8$、$r=2$ 时，整体只有 $1/0.6\approx1.67$ 倍加速；即使该部分耗时趋于零，总加速也不超过五倍。若额外引入格式转换开销，总收益还会更小。这个算例说明了层级内核加速与服务端到端加速为什么不能互换报告。

对语言模型应分别测量预填充延迟、首词元延迟、逐词元解码延迟和吞吐量，固定提示长度、生成长度与并发。权重稀疏不会自动减少所有注意力、通信和采样耗时。结构化缩宽有时能直接使用更小的稠密算子，但某些尺寸失去对齐优势也可能使利用率下降。最终应该测量导出的模型，而不仅是理论结构或带掩码的训练模型。

\subsection{可信计时与部署正确性}
计时前需要预热，以排除首次加载和编译等一次性成本；异步设备执行要在适当边界同步，否则记录的可能只是任务提交时间。多次测量应报告中位数或分位数，并说明是否包含数据搬运、分词、批处理与服务排队。不同精度、批量、算子后端和编译设置的结果不应直接归因于剪枝。

部署验证还应检查保存后重新加载的输出、掩码是否实际生效、稀疏格式是否被支持、是否出现静默退回稠密内核，以及模型配置是否和权重形状一致。对于不同数值内核允许合理浮点容差，但应以任务输出和数值误差共同判断。论文中的理论可加速性与当前系统里的实测收益是两类证据，应分别陈述。

\section{预算分配与可执行的算法流程}
\label{sec:prune-procedure}

\subsection{怎样决定各层剪多少}
设模型有 $L_b$ 个可剪层，第 $l$ 层原参数数为 $p_l$，保留数为 $k_l$，总预算为 $K_0$。一种近似做法是先估计各层在单独压缩时的损失曲线 $d_l(k_l)$，再求解 $\min\sum_l d_l(k_l)$ 且 $\sum_l k_l\leq K_0$。它把层之间的影响近似相加，便于离散动态规划或搜索，但忽略了前层变化对后层输入的影响。因此得到候选预算后，仍需联合剪枝并评估。

若每层有多个合法宽度选项，还可以为每个选项测量延迟，再在质量限制下选择组合。真实端到端延迟不一定等于各层独立延迟之和，因为算子融合、并行与缓存复用会发生变化。应把代理成本用于缩小搜索范围，再对少数候选执行真实测量。使用“重要性除以节省成本”的贪心比值是启发式，只有在某些可分离且满足额外结构的优化问题中才有更强保证。

\subsection{固定掩码的迭代剪枝流程}
一个完整训练后迭代流程的输入包括稠密模型、训练或恢复数据、校准与验证数据、目标稀疏率、候选参数集合、评分函数和恢复预算；输出应包括最终权重、掩码或结构配置以及测量记录。下面给出可直接转写为程序的步骤，假定每轮支持集只减少，不执行连接再生。
\begin{enumerate}
 \item 保存原模型、初始质量和资源测量，建立全部为一的候选掩码，并记录不可剪层、最小宽度与结构依赖。
 \item 在当前有效模型上收集所需统计量，依据评分规则和本轮累计预算选择保留集合；阈值处并列分数按固定索引规则处理。
 \item 更新掩码或结构，先测量即时损失；如出现非有限数值或超出预设质量界限，停止当前候选并检查尺度、层依赖和预算。
 \item 在固定掩码下恢复训练，持续施加约束，处理优化器状态，按预定验证准则选择检查点。
 \item 达到下一轮预算前重新评分，直至达到目标；随后导出部署表示，重新加载并评估质量、大小、内存和真实延迟。
\end{enumerate}
这套流程没有非凸全局收敛保证，停止条件应依据预算与验证表现制定，而不是假设剪得越多收益一定越好。即时损失和恢复后损失应分别记录，前者帮助比较评分是否预测了扰动，后者衡量最终可用性。

\subsection{无需标签的逐层重建流程}
若只能获得预训练模型与无标签输入，可使用局部重建。先确定需要压缩的线性层和合法稀疏模式，收集一批具有代表性的输入；按前向顺序为当前层记录输入激活，计算输入范数或二阶矩，选择掩码，并按算法决定是否补偿保留权重。处理完当前层后更新前缀，以规定的输入来源继续下一层。最后在独立数据上检验全模型行为。

使用式\eqref{eq:prune-wanda}时，主要统计成本约为 $O(n_cd)$，形成整层分数需要 $O(md)$；若显式形成 $d\times d$ 二阶矩，朴素成本约为 $O(n_cd^2)$，存储约为 $O(d^2)$，稠密分解通常为 $O(d^3)$。这些表达式只刻画基本运算规模，实际系统可分块、复用和分摊。不能因为不做梯度再训练，就把校准激活存储、矩阵分解和数据搬运的成本视为零。

\subsection{一个两层网络的完整手算方案}
设网络 $\mathbf h=\operatorname{ReLU}(\mathbf W_1\mathbf x+\mathbf b_1)$、$\widehat y=\mathbf w_2^\top\mathbf h+b_2$，输入维度为二，隐藏宽度为三。包含偏置时原参数数为 $3\times2+3+3+1=13$。若删除一个隐藏神经元并重建为宽度二的模型，参数数变为 $2\times2+2+2+1=9$，减少四个参数，而不是只少一个“神经元参数”。忽略偏置与激活操作，原乘加数为九，重建后为六。

为每个隐藏通道放置一个门，先在校准集上计算通道敏感度，并选出候选通道；分别测量删除前后的均方误差，再进行相同预算的恢复训练。导出时同步切片第一层行、第一层偏置和输出层列。若仅把输出层对应系数置零而仍计算三个隐藏神经元，函数可能已经不再依赖该神经元，但实际第一层算子仍未缩小。这个例子将评分、任务影响、参数减少和执行减少四件事明确分开。

\section{实验设计、理论边界与结果解释}
\label{sec:prune-evaluation}

\subsection{应当比较哪些基线}
至少保留原稠密模型、相同稀疏模式的随机剪枝和幅值剪枝，并为它们提供一致的恢复预算。结构化剪枝还应与直接训练同样小的架构比较，以区分收益来自继承权重、发现结构还是额外训练。关于重新训练小架构的研究提示了公平基线的重要性\cite{liu2019rethinking}；这一证据不能外推为所有预训练语言模型都适合从零训练，成本背景和任务规模必须同时考虑。

若方法使用更大的校准集、教师模型或更多恢复轮数，应把这些资源计入比较。一个分数复杂的方法仅在比基线多训练很多步后达到更好结果，尚不能证明评分本身更优。应把搜索成本、校准成本、恢复训练成本与最终推理成本分开，并给出主要超参数和随机种子。针对已有昂贵预训练模型的部署压缩，还应明确预训练成本是共同前提，还是被纳入整体算法成本。

\subsection{消融实验应对应具体机制}
若声称激活统计有价值，可保持掩码约束和恢复训练一致，比较纯幅值与激活加权；若声称误差补偿有效，可固定同一掩码，比较不补偿与补偿；若声称动态预算优于固定预算，可保持总保留数和评分一致，只改变层间分配。每次改变尽量对应一个机制，避免同时改变精度、数据、结构和训练步数，使结果无法归因。

对不同剪枝率绘制质量与成本曲线，通常比只报告一个点更有解释力。如果某模型质量不差且成本更低，可以说它支配另一个候选；所有未被支配的候选构成 Pareto 前沿。这里的成本可以是延迟、内存或能耗，但多种成本不能未经权重定义就合成一个“效率分数”。选择最终模型仍取决于使用者的质量下限与资源限制。

\subsection{有限样本上的保真不等于完整能力保真}
局部校准数据有限，模型可能在这些输入上重建准确，却在没有覆盖的方向上发生明显变化。设分类器对某个输入的正确类别 logit 比所有其他类别至少高出 $\gamma>0$，若剪枝使每个 logit 的绝对变化都小于 $\gamma/2$，预测类别保持不变。证明是任一正确与错误 logit 之差最多减少两倍扰动上界，仍大于零。这个条件强调了原分类间隔的重要性，也说明统一的输出均方误差不能直接保证所有样本类别不变。

更一般地，在输入范数有界、激活 Lipschitz、各层算子范数受控等条件下，可以逐层建立输出扰动上界，但最坏情况界常远大于实际误差。它们用于解释误差传播和提出约束，不宜当作实际性能的精确预测器。稀疏性减少了活跃参数数目，却不自动保证泛化改善；支持集搜索本身也使用了数据，且较小模型可能因容量不足而欠拟合。

\subsection{如何报告负结果和退化}
如果参数大幅减少而延迟不变，可能说明当前后端没有利用零模式，或瓶颈来自未被压缩的部分；这仍是有意义的系统结果。如果平均损失稳定而某个子任务明显退化，应报告该退化并检查校准覆盖、层预算和目标函数。不要把所有性能下降都归因于“剪得太多”，也不要用一次成功的校准子集选择掩盖不稳定性。

一份可复查的记录应包含数据划分、模型检查点、可剪参数范围、稀疏率分母、掩码粒度、恢复设置、导出格式、硬件与后端、质量指标和计时口径。多次运行的波动应反映数据抽样或训练随机性；对于完全确定的固定输入计时，重点则是系统噪声与分位数。本文中的数值算例用于验证机制和公式，没有将它们作为任何实际模型的实验结果。

\section{常见误区与练习解答}
\label{sec:prune-exercises}

\subsection{需要始终保持的几组区分}
零权重不等于删除了存储位置，稀疏支持集不等于更小的稠密架构，理论乘加减少不等于实测延迟同比减少。训练损失、校准重建误差、测试质量与部署成本属于不同评价层次。没有标签不等于没有数据，没有梯度再训练不等于不更新权重，保留原始初始化也不等于保留训练后的权重。只要在阅读方法时逐项定位这些对象，许多看似冲突的实验结论就会变得可以比较。

另一个常见误区是把评分公式当成完整算法。真正的算法还包括分数在哪些数据上计算、是否跨样本取绝对值、按哪一级预算排序、一次删多少、是否补偿、是否允许再生以及如何导出。相同公式搭配不同选择范围，可能得到差异很大的掩码；相同掩码搭配不同恢复训练，也可能得到不同质量。复现应保存这些决策，而不仅是公式名称。

\subsection{练习一：幅值与输入尺度}
设线性预测器为 $f(\mathbf X)=0.1X_1+X_2$，输入均值为零、不相关，方差分别为 $400$ 与 $1$，要求只保留一个连接且不补偿。请分别给出幅值剪枝和均方输出误差最小化的选择，并计算误差。

\begin{solve}
幅值法删除系数 $0.1$，均方输出误差为 $0.1^2\times400=4$。删除第二个系数的误差为 $1^2\times1=1$，所以按输出误差应删除第二个连接。无相关性时，修改第一个系数不能重建被删去的 $X_2$；这个例子再次说明参数空间距离与数据加权的函数距离不同。
\end{solve}

\subsection{练习二：正相关与联合删除}
设 $X_1,X_2$ 均值为零、方差均为一，相关系数为 $r\in(-1,1)$，模型为 $X_1+X_2$。删除第二个连接，允许第一个权重改为 $a$。求最优 $a$ 与最小均方输出误差，并解释 $r\to1$ 时的行为。

\begin{solve}
预测差为 $(a-1)X_1-X_2$，其均方值为 $(a-1)^2+1-2r(a-1)$。配方得 $(a-1-r)^2+1-r^2$，所以 $a^*=1+r$，最小误差为 $1-r^2$。当 $r\to1$，两个特征越来越相似，误差趋于零，最优权重趋于二。若同时删除两者，误差为 $\mathbb E[(X_1+X_2)^2]=2+2r$，不能把分别允许补偿的最小误差直接相加。
\end{solve}

\subsection{练习三：固定预算与多轮删除}
某层有一万个可剪权重，目标保留比例为 $0.25$。若计划四轮、每轮删除当前剩余权重的相同比例 $r$，忽略整数取整，求 $r$。若全模型还有两千个从不参与剪枝的参数，全模型最终稀疏率是多少？

\begin{solve}
由 $(1-r)^4=0.25$ 得 $r=1-0.25^{1/4}=1-1/\sqrt2\approx0.292893$。可剪部分保留两千五百个参数，连同受保护部分共四千五百个。相对全部一万两千个参数，最终稀疏率为 $1-4500/12000=62.5\%$，而可剪部分的稀疏率为 $75\%$。报告必须注明分母，实际逐轮取整还应保证最终累计预算准确。
\end{solve}

\subsection{练习四：通道删除与参数计数}
两个连续、无偏置的线性层形状分别为 $64\times32$ 和 $16\times64$，中间采用逐坐标激活。删除十六个隐藏神经元，外部输入输出维度不变。求前后参数量和单样本理想乘加次数，并说明为什么只给第一层行置零还不够。

\begin{solve}
原参数量为 $64\times32+16\times64=3072$，删除后为 $48\times32+16\times48=2304$，减少 $768$，即百分之二十五。忽略激活的乘加计数也从 $3072$ 降为 $2304$。若只在原稠密数组中置零，算子形状仍不变；必须联动切片第一层行和第二层列，并构建相应宽度的新算子，才得到这个更小的稠密计算图。
\end{solve}

\subsection{练习五：质量约束与可部署性}
假设两个候选模型在验证集准确率上分别下降 $0.2$ 与 $0.5$ 个百分点，实测端到端加速比分别为 $1.05$ 与 $1.6$，部署要求准确率下降不超过 $0.3$ 个百分点。是否应选择加速更多的候选？还需要检查什么？

\begin{solve}
第二个候选违反质量约束，所以在当前要求下不属于可行解。第一个候选虽然加速较小，但满足给定准确率条件；仍应检查延迟测量是否使用目标硬件与负载、子任务是否有额外质量限制，以及结果波动是否会使约束失效。如果两个候选都不能满足所有条件，应调整预算或方法，而不能只以加速比选择违反约束的模型。
\end{solve}

\subsection{从原理到独立实践}
开始实践时，可以先用一个小型网络实现固定掩码幅值剪枝，在相同预算下加入激活加权或一阶分数，最后比较补偿与恢复训练。每次都同时检查有效支持集、模型输出与部署成本，直到能够解释某个候选为什么更好或更差，再扩展到更复杂的二阶与结构搜索方法。学习的核心是建立一条可检验的因果链：某个统计量为何能预测删除影响，选择规则如何满足结构与预算，恢复如何改变结果，系统又如何把结构减少转化为资源收益。

\begin{thebibliography}{99}
\bibitem{lecun1989}
LeCun Y, Denker J S, Solla S A. Optimal Brain Damage. NIPS 1989, Advances in Neural Information Processing Systems 2.
\url{https://papers.nips.cc/paper/1989/hash/6c9882bbac1c7093bd25041881277658-Abstract.html}.

\bibitem{hassibi1992}
Hassibi B, Stork D G. Second Order Derivatives for Network Pruning: Optimal Brain Surgeon. NIPS 1992, Advances in Neural Information Processing Systems 5.
\url{https://proceedings.neurips.cc/paper/1992/hash/303ed4c69846ab36c2904d3ba8573050-Abstract.html}.

\bibitem{han2015}
Han S, Pool J, Tran J, Dally W J. Learning both Weights and Connections for Efficient Neural Networks. NeurIPS, 2015.
\url{https://arxiv.org/abs/1506.02626}.

\bibitem{wen2016}
Wen W, Wu C, Wang Y, Chen Y, Li H. Learning Structured Sparsity in Deep Neural Networks. NeurIPS, 2016.
\url{https://arxiv.org/abs/1608.03665}.

\bibitem{molchanov2017}
Molchanov P, et al. Pruning Convolutional Neural Networks for Resource Efficient Inference. ICLR, 2017.
\url{https://arxiv.org/abs/1611.06440}.

\bibitem{liu2017}
Liu Z, et al. Learning Efficient Convolutional Networks through Network Slimming. ICCV, 2017.
\url{https://arxiv.org/abs/1708.06519}.

\bibitem{louizos2018}
Louizos C, Welling M, Kingma D P. Learning Sparse Neural Networks through $L_0$ Regularization. ICLR, 2018.
\url{https://arxiv.org/abs/1712.01312}.

\bibitem{frankle2019}
Frankle J, Carbin M. The Lottery Ticket Hypothesis: Finding Sparse, Trainable Neural Networks. ICLR, 2019.
\url{https://arxiv.org/abs/1803.03635}.

\bibitem{lee2019}
Lee N, Ajanthan T, Torr P H S. SNIP: Single-shot Network Pruning based on Connection Sensitivity. ICLR, 2019.
\url{https://arxiv.org/abs/1810.02340}.

\bibitem{liu2019rethinking}
Liu Z, Sun M, Zhou T, Huang G, Darrell T. Rethinking the Value of Network Pruning. ICLR, 2019.
\url{https://arxiv.org/abs/1810.05270}.

\bibitem{wang2020}
Wang C, Zhang G, Grosse R. Picking Winning Tickets Before Training by Preserving Gradient Flow. ICLR, 2020.
\url{https://arxiv.org/abs/2002.07376}.

\bibitem{tanaka2020}
Tanaka H, Kunin D, Yamins D L K, Ganguli S. Pruning Neural Networks without Any Data by Iteratively Conserving Synaptic Flow. NeurIPS, 2020.
\url{https://arxiv.org/abs/2006.05467}.

\bibitem{evci2020}
Evci U, Gale T, Menick J, Castro P S, Elsen E. Rigging the Lottery: Making All Tickets Winners. ICML, 2020.
\url{https://proceedings.mlr.press/v119/evci20a.html}.

\bibitem{sanh2020}
Sanh V, Wolf T, Rush A M. Movement Pruning: Adaptive Sparsity by Fine-Tuning. NeurIPS, 2020.
\url{https://arxiv.org/abs/2005.07683}.

\bibitem{frantar2023}
Frantar E, Alistarh D. SparseGPT: Massive Language Models Can Be Accurately Pruned in One-Shot. ICML, 2023.
\url{https://proceedings.mlr.press/v202/frantar23a.html}.

\bibitem{sun2024}
Sun M, Liu Z, Bair A, Kolter J Z. A Simple and Effective Pruning Approach for Large Language Models. ICLR, 2024.
\url{https://arxiv.org/abs/2306.11695}.

\bibitem{pytorchprune}
PyTorch Contributors. Pruning Tutorial. 官方文档，访问日期：2026 年 10 月 3 日。
\url{https://docs.pytorch.org/tutorials/intermediate/pruning_tutorial.html}.
\end{thebibliography}

\end{document}
